\documentclass{article}

\usepackage[english]{babel}
\usepackage[utf8]{inputenc}
\usepackage{polski}
\usepackage[T1]{fontenc}
 
\usepackage[margin=1.5in]{geometry} 

\usepackage{color} 
\usepackage{amsmath}
\numberwithin{equation}{subsection}

\usepackage{amssymb}
\usepackage{amsfonts}                                                                   
\usepackage{graphicx}                                                             
\usepackage{booktabs}
\usepackage{amsthm}
\usepackage{pdfpages}
\usepackage{wrapfig}
\usepackage{hyperref}
\usepackage{etoolbox}
\usepackage{tikz}

\makeatletter
\newenvironment{definition}[1]{%
    \trivlist
    \item[\hskip\labelsep\textbf{Definition. #1.}]
    \ignorespaces
}{%
    \endtrivlist
}
\newenvironment{fact}[1]{%
    \trivlist
    \item[\hskip\labelsep\textbf{Fact. #1.}]
    \ignorespaces
}{%
    \endtrivlist
}
\newenvironment{theorem}[1]{%
    \trivlist
    \item[\hskip\labelsep\textbf{Theorem. #1.}]
    \ignorespaces
}{%
    \endtrivlist
}
\newenvironment{information}[1]{%
    \trivlist
    \item[\hskip\labelsep\textbf{Information. #1.}]
    \ignorespaces
}{%
    \endtrivlist
}
\newenvironment{identities}[1]{%
    \trivlist
    \item[\hskip\labelsep\textbf{Identities. #1.}]
    \ignorespaces
}{%
    \endtrivlist
}
\makeatother

\title{Teoretyczne Podstawy Informatyki}  
\author{Rafał Włodarczyk}
\date{INA 6, 2026}

\begin{document}

\maketitle

\tableofcontents

\vspace{1cm}

\noindent
Uniwersalne pojęcie \textbf{obliczenia} nie istnieje. Musimy wpierw wybrać model - Maszyny Turinga, Rachunku Lambda.

\section{Maszyna Turinga}

Maszyna Turinga to $(Q,q_0,\Sigma,\delta)$:
\begin{enumerate}
    \item $Q$ - skończony zbiór stanów
    \item $q_0 \in Q$ - stan początkowy
    \item $\Sigma$ - skończony alfabet (oraz $\smile$ (brak) dla pustych komórek)
    \item $\delta : Q \times (\Sigma \cup \{\smile\}) \to (Q \cup \{\text{tak},\text{nie},\text{h}\}) \times (\Sigma \cup \{\smile\}) \times \{\to, \leftarrow, -\}$
\end{enumerate}
oraz mamy nieskończoną taśmę.

\subsection{Palindrom nad Alfabetem}

Przykład. Palindrom nad Alfabetem $\{0,1\}$\\
\begin{tabular}{|c|ccc|}
\hline
stan & 0 & 1 & $\smile$ \\
\hline
$p$ & $(q_0,\smile,\to)$ & $(q_1,\smile,\to)$ & $\{\text{tak}\}$ \\
$q_0$ & $(q_0, 0, \to)$ & $(q_0, 1, \to)$ & $(q_2, \smile, \leftarrow)$ \\
$q_1$ & $(q_0, 0, \to)$ & $(q_1, 1, \to)$ & $(q_3, \smile, \leftarrow)$ \\
$q_2$ & $(r,\smile, \leftarrow)$ & $\{\text{nie}\}$ & $\{\text{tak}\}$ \\ 
$q_3$ & $\{\text{nie}\}$ & $(r,\smile,\leftarrow)$ & $\{\text{tak}\}$ \\ 
$r$ & $(r,0,\leftarrow)$ & $(r, 1, \leftarrow)$ & $(p,\smile,\to)$ \\
\hline
\end{tabular}
\\\\
Uproszczenie zapisu: $\{\text{tak}\}\ \equiv (\text{tak}, \smile, -)$. Czas rozpoznania palindromu: $\mathcal{O}(n^2)$

\subsection{Wynik działania Maszyny Turinga}

Dla TM $M$ i ciągu wejściowego $X$ przez $M(X)$ oznaczamy wynik działania $M$ na $x$.
\begin{itemize}
    \item $M(x) = \text{tak}$ - mówimy, że $M$ akceptuje $x$
    \item $M(x) = \text{nie}$ - mówimy, że $M$ odrzuca $x$
    \item $M(x) = \text{y}$ - mówimy, że $M$ zatrzymała się na $h$ i na taśmie jest $y$ (bez zbędnych $\smile$ z lewej i prawej strony). 
    Mówimy, że $M(x)$ obliczyło $y$.
    \item $M(x) = \nearrow$ - jeśli $M$ na $x$ nie kończy nigdy pracy (zapętla się).
\end{itemize}

\subsection{Dodawanie Binarne i napisy nad 0,1}

Przykład. Dodawanie binarne.\\
\begin{tabular}{|c|ccc|}
\hline
stan & 0 & 1 & $\smile$ \\
\hline
$q_0$ & $(q_0, 0, \to)$ & $(q_0,1,\to)$ & $(q_1,\smile, \leftarrow)$ \\
$q_1$ & $(q_2, 1, \leftarrow)$ & $(q_1, 0, \leftarrow)$ & $(h, 1, -)$ \\ 
$q_2$ & $(q_2, 0, \leftarrow)$ & $(q_2, 1, \leftarrow)$ & $(h, \smile, \to)$ \\
\hline
\end{tabular}
\\\\
Dla maszyny $M$ zachodzi $M(101)=110$. Jeśli zmienie $h$ na $q_0$ to zapętlamy dodawanie.
\\\\
Przykład. Generujemy napisy nad $\{0,1\}$\\
\begin{tabular}{|c|ccc|}
\hline
stan & 0 & 1 & $\smile$ \\
\hline
$q_0$ & $(q_0, 0, \to)$ & $(q_0,1,\to)$ & $(q_1,\smile, \leftarrow)$ \\
$q_1$ & $(q_2, 1, \leftarrow)$ & $(q_1, 0, \leftarrow)$ & $(h, \color{red}{0}\color{black}, -)$ \\ 
$q_2$ & $(q_2, 0, \leftarrow)$ & $(q_2, 1, \leftarrow)$ & $(h, \smile, \to)$ \\
\hline
\end{tabular}
\\\\
$M$: $0,1,00,01,10,11,000,\dots$

\subsection{Konfiguracja Maszyny Turinga}

Konfiguracją TM $M$ nazywamy trójkę $(q,\alpha, \beta)$, taką że:
\begin{itemize}
    \item $q$ jest aktualnym stanem
    \item $\alpha$ jest zawartością taśmy przed głowicą
    \item $\beta$ jest zawartością tasmy po głowicy
    \item $\vdash_M$ oznacza relację przejścia w jednym kroku przez maszynę $M$ między konfiguracjami
\end{itemize}
Analogicznie $\vdash_M^i$ to wyprowadzenie w $i$-krokach, $\vdash_M^{*}$ w dowolnej ilości kroków.
\begin{align*}
    x = 010& \\
    (p,\varepsilon, 010) &\vdash \\
    (q_0, \varepsilon, 10) &\vdash \\
    (q_0, 1, 0) &\vdash \\
    (q_0, 10, \varepsilon) &\vdash \\
    (q_2, 1, 0) &\vdash \\
    (r, \varepsilon, 1) &\vdash \dots
\end{align*}
$M(x) = \text{tak} \iff (q_0,\varepsilon, x) \vdash_M^{*} (\text{tak}, *, *)$\\
$M(x) = y \iff (q_0, \varepsilon, x) \vdash_M^{*} (h, \alpha, \beta) \land y=\alpha\beta$

\section{k-Taśmowa Maszyna Turinga}

Zmienia się $\delta$ na:
\begin{align*}
    \delta : Q \times \left(\sum \cup \{\smile\}\right)^k 
    \to
    \left(Q\sum \{\text{tak, nie, h}\}\right) \times
    \left(\left(\sum \cup \{\smile\}\right) 
    \times \{\to,\leftarrow, -\}\right)^k
\end{align*}
Przykład. Palindrom nad alfabetem $\{0,1\}$\\
\begin{tabular}{cccccccccc}
\hline
\text{stan} & 00 & 01 & 0$\smile$ & 10 & 11 \\
$q_0$ & & & $q_0,\to^0,\to^0$ & & & $q_0,\to^1,\to^1$ \\
$q_1$ & $q_1, \leftarrow^0, -^0$ & $\leftarrow^0 -^1$ & & $q_1, \leftarrow^1 -^0$ & $q_1, \leftarrow^1 -^1$ \\
$q_2$ & $q_2, \to^{\{\smile\}}, \leftarrow^{\{\smile\}}$ & nie & - & nie &  $q_2, \to^{\{\smile\}}\leftarrow^{\{\smile\}}$ \\
\hline
\text{stan} & 1$\smile$ & $\smile$0 & $\smile$1 & $\smile\smile$ \\
$q_0$ & $q_0,\leftarrow^{\{\smile\}},\leftarrow^{\{\smile\}}$\\ 
$q_1$ & $q_2, \to^{\{\smile\}}, -^0$ &  $q_2, \to^{\{\smile\}}, -^1$ & tak \\
$q_2$ & - & - & - & tak
\end{tabular}

\subsection{Zastąpienie k-Taśmowej przez jednotaśmową}

Każda $k$-taśmowa TM $M$ która na $x$ pracuje w czasie $f(|x|)$ można zastąpić $1$-taśmową TM $M'$ taką, że:
\begin{enumerate}
    \item $\forall x M(x) = M'(x)$
    \item $M'(x)$ wykona co najwyżej $O(f(|x|)^2)$ kroków.
\end{enumerate}

\section{Niedeterminizm}

\begin{align*}
    \delta : Q \times \left(\sum \cup \{\smile\}\right) \to
    2^{(Q\cup\{\text{tak, nie, h}\}) \times (\Sigma \cup \{\smile\}) \times \{\to,\leftarrow,-\}}
\end{align*}
NTM $M$ akceptuje $x$ jeśli istnieje obliczenie $M(x)$ kończące się w stanie \textit{tak},
odrzuca jeśli każde obliczenie kończy się w stanie \textit{nie}.




\end{document}
